\chapter*{ПРИЛОЖЕНИЕ А}
\addcontentsline{toc}{chapter}{ПРИЛОЖЕНИЕ А}

\section*{Презентация к курсовой работе}
Презентация содержит 12 слайдов.

\chapter*{ПРИЛОЖЕНИЕ B}
\addcontentsline{toc}{chapter}{ПРИЛОЖЕНИЕ B}

\section*{Листинг к курсовой работе}

\begin{lstlisting}[label=lst:cnst1,caption=Запросы для создания ограничений таблиц (часть 1)]
ALTER TABLE client
    ALTER COLUMN client_login SET NOT NULL,
    ALTER COLUMN client_password SET NOT NULL,
    ALTER COLUMN mail SET NOT NULL,
    ALTER COLUMN privilege SET DEFAULT 'UNSIGNED',
    ADD CONSTRAINT uk_client_login
        UNIQUE (client_login),
    ADD CONSTRAINT uk_client_mail
        UNIQUE (mail);

ALTER TABLE chart
    ALTER COLUMN city SET NOT NULL,
    ALTER COLUMN title SET NOT NULL,
    ALTER COLUMN svg_content SET DEFAULT '<svg width="200" height="200"><text y="16">Пустая схема</text></svg>',
    ADD CONSTRAINT ck_valid_svg
        CHECK (svg_content::xml IS DOCUMENT),
    ADD CONSTRAINT uk_chart_city_title
        UNIQUE (city, title);

ALTER TABLE branch
    ALTER COLUMN title SET NOT NULL,
    ALTER COLUMN color SET DEFAULT 0,
    ALTER COLUMN access SET DEFAULT 'ACCESSIBLE';

ALTER TABLE station
    ALTER COLUMN duty_id DROP NOT NULL,
    ALTER COLUMN title SET NOT NULL,
    ALTER COLUMN occupancy SET DEFAULT 5,
    ALTER COLUMN access SET DEFAULT 'ACCESSIBLE',
    ALTER COLUMN open_time SET NOT NULL,
    ALTER COLUMN close_time SET NOT NULL,
    ADD CONSTRAINT ck_station_different_open_close_time
        CHECK (open_time <> close_time),
    ADD CONSTRAINT fk_station_duty_id
        FOREIGN KEY (duty_id) REFERENCES client (id) ON DELETE SET NULL;

ALTER TABLE transition
    ALTER COLUMN duty_id DROP NOT NULL,
    ALTER COLUMN occupancy SET DEFAULT 5,
    ALTER COLUMN access SET DEFAULT 'ACCESSIBLE',
    ALTER COLUMN duration SET NOT NULL,
    ALTER COLUMN open_time SET NOT NULL,
    ALTER COLUMN close_time SET NOT NULL,
    ADD CONSTRAINT ck_transition_different_open_close_time
        CHECK (open_time <> close_time),
    ADD CONSTRAINT fk_transition_duty_id
        FOREIGN KEY (duty_id) REFERENCES client (id) ON DELETE SET NULL;
\end{lstlisting}

\clearpage

\begin{lstlisting}[label=lst:cnst2,caption=Запросы для создания ограничений таблиц (часть 2)]
ALTER TABLE way
    ALTER COLUMN client_id SET NOT NULL,
    ALTER COLUMN chart_id SET NOT NULL,
    ALTER COLUMN title SET NOT NULL,
    ALTER COLUMN init_date SET NOT NULL,
    ADD CONSTRAINT uk_way_client_id_title
        UNIQUE (client_id, title),
    ADD CONSTRAINT fk_way_client_id
        FOREIGN KEY (client_id) REFERENCES client (id) ON DELETE CASCADE,
    ADD CONSTRAINT fk_way_chart_id
        FOREIGN KEY (chart_id) REFERENCES chart (id) ON DELETE CASCADE;


ALTER TABLE way_item
    ALTER COLUMN way_id SET NOT NULL,
    ALTER COLUMN nexus SET NOT NULL,
    ALTER COLUMN step_nomer SET NOT NULL,
    ADD CONSTRAINT uk_way_item_way_id_step_nomer
        UNIQUE (way_id, step_nomer),
    ADD CONSTRAINT fk_way_item_way_id
        FOREIGN KEY (way_id) REFERENCES way (id) ON DELETE CASCADE;


ALTER TABLE chart_branch
    ALTER COLUMN chart_id SET NOT NULL,
    ALTER COLUMN branch_id SET NOT NULL,
    ADD CONSTRAINT uk_chart_branch_chart_id_branch_id
        UNIQUE (chart_id, branch_id), -- Ни одна схема не имеет дубликатов веток
    ADD CONSTRAINT uk_chart_branch_branch_id
        UNIQUE (branch_id), -- Каждая ветка связывается со схемой только один раз
    ADD CONSTRAINT fk_chart_branch_chart_id
        FOREIGN KEY (chart_id) REFERENCES chart (id) ON DELETE CASCADE,
    ADD CONSTRAINT fk_chart_branch_branch_id
        FOREIGN KEY (branch_id) REFERENCES branch (id) ON DELETE CASCADE;


ALTER TABLE branch_station
    ALTER COLUMN branch_id SET NOT NULL,
    ALTER COLUMN station_id SET NOT NULL,
    ADD CONSTRAINT uk_branch_station_branch_id_station_id
        UNIQUE (branch_id, station_id), -- Ни одна ветка не имеет дубликатов станций
    ADD CONSTRAINT uk_branch_station_station_id
        UNIQUE (station_id), -- Каждая станция связывается с веткой только один раз
    ADD CONSTRAINT fk_branch_station_branch_id
        FOREIGN KEY (branch_id) REFERENCES branch (id) ON DELETE CASCADE,
    ADD CONSTRAINT fk_branch_station_station_id
        FOREIGN KEY (station_id) REFERENCES station (id) ON DELETE CASCADE;


ALTER TABLE railway
    ALTER COLUMN from_id SET NOT NULL,
    ALTER COLUMN to_id SET NOT NULL,
    ALTER COLUMN duration SET NOT NULL,
    ADD CONSTRAINT uk_railway_from_id_to_id
        UNIQUE (from_id, to_id),
    ADD CONSTRAINT ck_railway_different_stations
        CHECK (from_id <> to_id),
    ADD CONSTRAINT fk_railway_from_id
        FOREIGN KEY (from_id) REFERENCES station (id) ON DELETE CASCADE,
    ADD CONSTRAINT fk_railway_to_id
        FOREIGN KEY (to_id) REFERENCES station (id) ON DELETE CASCADE;
\end{lstlisting}

\clearpage

\begin{lstlisting}[label=lst:cnst3,caption=Запросы для создания ограничений таблиц (часть 3)]
ALTER TABLE station_transition
    ALTER COLUMN station_id SET NOT NULL,
    ALTER COLUMN transition_id SET NOT NULL,
    ADD CONSTRAINT uk_station_transition_station_id_transition_id
        UNIQUE (station_id, transition_id),
    ADD CONSTRAINT fk_station_transition_station_id
        FOREIGN KEY (station_id) REFERENCES station (id) ON DELETE CASCADE,
    ADD CONSTRAINT fk_station_transition_transition_id
        FOREIGN KEY (transition_id) REFERENCES transition (id) ON DELETE CASCADE;


ALTER TABLE way_item_station
    ALTER COLUMN way_item_id SET NOT NULL,
    ALTER COLUMN station_id SET NOT NULL,
    ADD CONSTRAINT uk_way_item_station_way_item_id_station_id
        UNIQUE (way_item_id, station_id),
    ADD CONSTRAINT uk_way_item_station_way_item_id
        UNIQUE (way_item_id),
    ADD CONSTRAINT fk_way_item_station_way_item_id
        FOREIGN KEY (way_item_id) REFERENCES way_item (id) ON DELETE CASCADE,
    ADD CONSTRAINT fk_way_item_station_station_id
        FOREIGN KEY (station_id) REFERENCES station (id) ON DELETE CASCADE;


ALTER TABLE way_item_railway
    ALTER COLUMN way_item_id SET NOT NULL,
    ALTER COLUMN railway_id SET NOT NULL,
    ADD CONSTRAINT uk_way_item_railway_way_item_id_railway_id
        UNIQUE (way_item_id, railway_id),
    ADD CONSTRAINT uk_way_item_railway_way_item_id
        UNIQUE (way_item_id),
    ADD CONSTRAINT fk_way_item_railway_way_item_id
        FOREIGN KEY (way_item_id) REFERENCES way_item (id) ON DELETE CASCADE,
    ADD CONSTRAINT fk_way_item_railway_railway_id
        FOREIGN KEY (railway_id) REFERENCES railway (id) ON DELETE CASCADE;


ALTER TABLE way_item_transition
    ALTER COLUMN way_item_id SET NOT NULL,
    ALTER COLUMN transition_id SET NOT NULL,
    ADD CONSTRAINT uk_way_item_transition_way_item_id_transition_id
        UNIQUE (way_item_id, transition_id),
    ADD CONSTRAINT uk_way_item_transition_way_item_id
        UNIQUE (way_item_id),
    ADD CONSTRAINT fk_way_item_transition_way_item_id
        FOREIGN KEY (way_item_id) REFERENCES way_item (id) ON DELETE CASCADE,
    ADD CONSTRAINT fk_way_item_transition_transition_id
        FOREIGN KEY (transition_id) REFERENCES transition (id) ON DELETE CASCADE;
\end{lstlisting}

\clearpage

\begin{lstlisting}[label=lst:add_chart_json:1,caption=Реализация функции добавления схемы (часть 1)]
CREATE OR REPLACE FUNCTION add_chart_json(
    p_json_data TEXT
) RETURNS INT AS $$
DECLARE
    v_json_data JSONB := p_json_data::JSONB;

    v_chart_id      INT;   -- айди созданной записи схемы
    v_branch_id     INT;   -- айди созданной записи ветки
    v_station_id    INT;   -- айди созданной записи станции
    v_transition_id INT;   -- айди созданной записи перехода

    v_branch_item     RECORD;
    v_station_item    RECORD;
    v_railway_item    RECORD;
    v_transition_item RECORD;
BEGIN
    IF p_json_data IS NULL OR NOT (p_json_data ~ '^[\s]*\{.*\}[\s]*$') THEN
        RAISE EXCEPTION 'Некорректный JSON: входные данные пусты или не объект';
    END IF;

    INSERT INTO chart(city, title)
    VALUES(
        v_json_data->>'city',
        v_json_data->>'title'
    )
    RETURNING id INTO v_chart_id;

    FOR v_branch_item IN
        SELECT * FROM jsonb_array_elements(v_json_data->'branches')
    LOOP

        INSERT INTO branch(title, color, access)
        VALUES(
            (v_branch_item.value)->>'title',
            hex_color_to_int((v_branch_item.value)->>'color'),
            ((v_branch_item.value)->>'accesstype')::access_type
        )
        RETURNING id INTO v_branch_id;

        INSERT INTO chart_branch(chart_id, branch_id)
        VALUES(
            v_chart_id,
            v_branch_id
        );

        FOR v_station_item IN
            SELECT * FROM jsonb_array_elements((v_branch_item.value)->'stations')
        LOOP

            INSERT INTO station(title, occupancy, access, open_time, close_time)
            VALUES(
                (v_station_item.value)->>'title',
                ((v_station_item.value)->>'occupancy')::SMALLINT,
                ((v_station_item.value)->>'accesstype')::access_type,
                ((v_station_item.value)->>'opentime')::TIME,
                ((v_station_item.value)->>'closetime')::TIME
            )
            RETURNING id INTO v_station_id;
\end{lstlisting}

\clearpage

\begin{lstlisting}[label=lst:add_chart_json:2,caption=Реализация функции добавления схемы (часть 2),firstnumber=59]
            INSERT INTO branch_station(branch_id, station_id)
            VALUES(
                v_branch_id,
                v_station_id
            );

        END LOOP;

        FOR v_railway_item IN
            SELECT * FROM jsonb_array_elements((v_branch_item.value)->'railways')
        LOOP

            INSERT INTO railway(from_id, to_id, duration)
            VALUES(
                get_station_id_by_branch_id_station_title(v_branch_id, (v_railway_item.value)->>'from'),
                get_station_id_by_branch_id_station_title(v_branch_id, (v_railway_item.value)->>'to'),
                ((v_railway_item.value)->>'duration')::INTERVAL
            );

        END LOOP;

    END LOOP;

    FOR v_transition_item IN
        SELECT * FROM jsonb_array_elements(v_json_data->'transitions')
    LOOP

        INSERT INTO transition(occupancy, access, duration, open_time, close_time)
        VALUES(
            ((v_transition_item.value)->>'occupancy')::SMALLINT,
            ((v_transition_item.value)->>'accesstype')::access_type,
            ((v_transition_item.value)->>'duration')::INTERVAL,
            ((v_transition_item.value)->>'opentime')::TIME,
            ((v_transition_item.value)->>'closetime')::TIME
        )
        RETURNING id INTO v_transition_id;

        INSERT INTO station_transition(station_id, transition_id)
        VALUES (
            get_station_id_by_branch_title_station_title((v_transition_item.value)
                ->>'branchsrc', (v_transition_item.value)->>'stationsrc'),
            v_transition_id
        );

        INSERT INTO station_transition(station_id, transition_id)
        VALUES (
            get_station_id_by_branch_title_station_title((v_transition_item.value)
                ->>'branchdst', (v_transition_item.value)->>'stationdst'),
            v_transition_id
        );

    END LOOP;

    RETURN v_chart_id;
END;
$$ LANGUAGE plpgsql;
\end{lstlisting}

\clearpage

\begin{lstlisting}[label=lst:add_route_json:1,caption=Реализация функции добавления маршрута (часть 1),firstnumber=1]
CREATE OR REPLACE FUNCTION add_route_json(
    p_client_id INT,
    p_chart_id INT,
    p_json_data TEXT
) RETURNS INT AS $$
DECLARE
    v_json_data JSONB := p_json_data::JSONB;

    v_way_id INT;          -- айди созданной записи маршрута
    v_item RECORD;         -- обход массива звеньев маршрута  
    v_way_item_id INT;     -- айди созданного звена
    v_step_nomer INT := 1; -- номер звена маршрута
    v_branch_id INT;       -- айди найденной ветки
    v_station_id INT;      -- айди найденной станции
    v_station_from_id INT; -- айди станции: <ОТ>
    v_station_to_id INT;   -- айди станции: <НА>
    v_railway_id INT;      -- айди найденного переезда
    v_transition_id INT;   -- айди найденного перехода
BEGIN
    -- Проверка конвертации
    IF p_json_data IS NULL OR NOT (p_json_data ~ '^[\s]*\{.*\}[\s]*$') THEN
        RAISE EXCEPTION 'Некорректный JSON: входные данные пусты или не объект';
    END IF;

    -- Создать запись нового маршрута
    INSERT INTO way (client_id, chart_id, title, duration, init_date)
    VALUES (
        p_client_id,
        p_chart_id,
        v_json_data->>'Title',
        (v_json_data->>'Duration')::INTERVAL,
        date_trunc('second', NOW())::TIMESTAMPTZ
    )
    RETURNING id INTO v_way_id;

    -- Вставка звеньев маршрута
    FOR v_item IN SELECT * FROM jsonb_array_elements(v_json_data->'RouteItems')
    LOOP
        -- Создать запись о звенье маршрута
        INSERT INTO way_item (way_id, nexus, step_nomer)
        VALUES (
            v_way_id,
            CASE 
                WHEN (v_item.value)->>'$type' = 'station' THEN 'STATION'::nexus_type
                WHEN (v_item.value)->>'$type' = 'railway' THEN 'RAILWAY'::nexus_type
                WHEN (v_item.value)->>'$type' = 'transition' THEN 'TRANSITION'::nexus_type
            END,
            v_step_nomer
        )
        RETURNING id INTO v_way_item_id;

        -- Конкретизация и связывание звяна с реальными станциями, переездами и перехода схемы
        CASE
            -- Станция
            WHEN (v_item.value)->>'$type' = 'station' THEN
                -- Поиск связываемой ветки (обновлять айди нужно только на каждую станцию и переход)
                SELECT
                    b.id INTO v_branch_id
                FROM 
                    branch AS b
                WHERE
\end{lstlisting}

\clearpage

\begin{lstlisting}[label=lst:add_route_json:2,caption=Реализация функции добавления маршрута (часть 2),firstnumber=62]
                    b.title = (v_item.value)->>'BranchTitle' AND
                    EXISTS(
                        SELECT
                            1
                        FROM
                            chart_branch AS cb
                        WHERE
                            cb.chart_id = p_chart_id AND cb.branch_id = b.id
                    );

                -- Поиск связываемой станции
                SELECT
                    s.id INTO v_station_id
                FROM
                    station AS s
                INNER JOIN
                    branch_station AS bs ON s.id = bs.station_id
                WHERE
                    bs.branch_id = v_branch_id AND
                    s.title = (v_item.value)->>'Title';

                -- Вставка звена маршрута <Станция>
                INSERT INTO way_item_station (way_item_id, station_id)
                VALUES (v_way_item_id, v_station_id);

            -- Переезд
            WHEN (v_item.value)->>'$type' = 'railway' THEN
                -- Поиск отправной станции
                SELECT
                    s.id INTO v_station_from_id
                FROM
                    station AS s
                INNER JOIN
                    branch_station AS bs ON s.id = bs.station_id
                WHERE
                    bs.branch_id = v_branch_id AND
                    s.title = (v_item.value)->>'FromStationTitle';
                
                -- Поиск станции назначения
                SELECT
                    s.id INTO v_station_to_id
                FROM
                    station AS s
                INNER JOIN
                    branch_station AS bs ON s.id = bs.station_id
                WHERE
                    bs.branch_id = v_branch_id AND
                    s.title = (v_item.value)->>'ToStationTitle';

                -- Поиск связываемого переезда
                SELECT
                    r.id INTO v_railway_id
                FROM
                    railway AS r
                WHERE
                    (r.from_id = v_station_from_id AND r.to_id = v_station_to_id) OR
                    (r.from_id = v_station_to_id AND r.to_id = v_station_from_id);

                IF v_railway_id IS NULL THEN
                    RAISE EXCEPTION 'На шаге % не найден переезд между станциями % и % на ветке %',
                        v_step_nomer, v_station_from_id, v_station_to_id, v_branch_id;
                END IF;
\end{lstlisting}

\clearpage

\begin{lstlisting}[label=lst:add_route_json:3,caption=Реализация функции добавления маршрута (часть 3),firstnumber=124]
                -- Вставка звена маршрута <Переезд>
                INSERT INTO way_item_railway (way_item_id, railway_id)
                VALUES (v_way_item_id, v_railway_id);
            
            -- Переход
            WHEN (v_item.value)->>'$type' = 'transition' THEN               
                -- Поиск первой связанной станции
                SELECT
                    s.id INTO v_station_from_id
                FROM
                    station AS s
                WHERE
                    s.Title = (v_item.value)->>'FromStationTitle' AND
                    EXISTS(
                        SELECT
                            1
                        FROM
                            branch AS b
                        WHERE
                            b.title = (v_item.value)->>'FromBranchTitle' AND
                            EXISTS (
                                SELECT
                                    1
                                FROM
                                    branch_station AS bs
                                WHERE
                                    bs.branch_id = b.id AND bs.station_id = s.id
                            )
                    );
                
                -- Поиск второй связанной станции
                SELECT
                    s.id INTO v_station_to_id
                FROM
                    station AS s
                WHERE
                    s.Title = (v_item.value)->>'ToStationTitle' AND
                    EXISTS(
                        SELECT
                            1
                        FROM
                            branch AS b
                        WHERE
                            b.title = (v_item.value)->>'ToBranchTitle' AND
                            EXISTS (
                                SELECT
                                    1
                                FROM
                                    branch_station AS bs
                                WHERE
                                    bs.branch_id = b.id AND bs.station_id = s.id
                            )
                    );
                
                -- Поиск связываемого перехода
                SELECT
                    st.transition_id INTO v_transition_id
                FROM
                    station_transition AS st
                WHERE
                    st.station_id = v_station_from_id OR st.station_id = v_station_to_id
                GROUP BY
                    st.transition_id
\end{lstlisting}

\clearpage

\begin{lstlisting}[label=lst:add_route_json:4,caption=Реализация функции добавления маршрута (часть 4),firstnumber=187]
                HAVING
                    COUNT(*) = 2; -- Только у подходящего перехода количесво связей будет равно 2

                -- Вставка звена маршрута <Переход>
                INSERT INTO way_item_transition (way_item_id, transition_id)
                VALUES (v_way_item_id, v_transition_id);

        END CASE;

        v_step_nomer := v_step_nomer + 1;
    END LOOP;

    RETURN v_way_id;
END;
$$ LANGUAGE plpgsql;
\end{lstlisting}

\begin{lstlisting}[label=lst:get_chart_json_by_id:1,caption=Реализация функции получения схемы (часть 1),firstnumber=1]
CREATE OR REPLACE FUNCTION get_chart_json_by_id(chart_id INT)
RETURNS JSONB AS $$
DECLARE
    result_json JSONB;
BEGIN
    SELECT
        jsonb_build_object(
            'city', c.city,
            'title', c.title,
            'branches', COALESCE((
            SELECT
                jsonb_agg(
                    jsonb_build_object(
                        'title', b.title,
                        'color', UPPER((SELECT int_color_to_hex(b.color))),
                        'accesstype', b.access,
                        'stations', COALESCE((
                        SELECT
                            jsonb_agg(
                                jsonb_build_object(
                                    'title', s.title,
                                    'occupancy', s.occupancy,
                                    'accesstype', s.access,
                                    'opentime', (SELECT TO_CHAR(s.open_time, 'HH24:MI')),
                                    'closetime', (SELECT TO_CHAR(s.close_time, 'HH24:MI'))
                                )
                            )
                        FROM
                            station AS s
                        WHERE
                            EXISTS (
                                SELECT
                                    1
                                FROM
                                    branch_station AS bs
                                WHERE
                                    bs.station_id = s.id AND bs.branch_id = b.id
                            )
                        ), '[]'::JSONB),
                        'railways', COALESCE((
                        SELECT
                            jsonb_agg(
                                jsonb_build_object(
\end{lstlisting}

\clearpage

\begin{lstlisting}[label=lst:get_chart_json_by_id:2,caption=Реализация функции получения схемы (часть 2),firstnumber=44]
                                    'from', (
                                    SELECT
                                        s.title
                                    FROM
                                        station AS s
                                    WHERE
                                        s.id = r.from_id
                                    ),
                                    'to', (
                                    SELECT
                                        s.title
                                    FROM
                                        station AS s
                                    WHERE
                                        s.id = r.to_id
                                    ),
                                    'duration', r.duration
                                )
                            )
                        FROM
                            railway AS r
                        WHERE
                            EXISTS (
                                SELECT
                                    1
                                FROM
                                    branch_station AS bs
                                JOIN
                                    station AS s ON b.id = bs.branch_id AND s.id = bs.station_id AND s.id = r.from_id
                            )
                        ), '[]'::JSONB)
                    )
                )
            FROM
                branch AS b
            JOIN
                chart_branch AS cb ON cb.chart_id = c.id AND cb.branch_id = b.id
            ), '[]'::JSONB),
            'transitions', COALESCE((
            WITH
                adj AS (
                    SELECT
                        bs.branch_id,
                        bs.station_id,
                        t.id AS transition_id,
                        t.occupancy,
                        t.access,
                        t.duration,
                        t.open_time,
                        t.close_time
                    FROM
                        chart_branch AS cb
                    JOIN
                        branch_station AS bs ON cb.chart_id = c.id AND cb.branch_id = bs.branch_id
                    JOIN
                        station_transition AS st ON st.station_id = bs.station_id
                    JOIN
                        transition AS t ON t.id = st.transition_id
                ),
\end{lstlisting}

\clearpage

\begin{lstlisting}[label=lst:get_chart_json_by_id:3,caption=Реализация функции получения схемы (часть 3),firstnumber=103]
                r_adj AS (
                    SELECT
                        adj.branch_id,
                        adj.station_id,
                        adj.transition_id,
                        rank() OVER (PARTITION BY transition_id ORDER BY branch_id, station_id)
                    FROM
                        adj
                ),
                d_adj AS (
                    SELECT DISTINCT ON (adj.transition_id)
                        *
                    FROM
                        adj
                    ORDER BY
                        adj.transition_id
                )
            SELECT
                jsonb_agg(
                    jsonb_build_object(
                        'occupancy', d_adj.occupancy,
                        'accesstype', d_adj.access,
                        'duration', d_adj.duration,
                        'opentime', (SELECT TO_CHAR(d_adj.open_time, 'HH24:MI')),
                        'closetime', (SELECT TO_CHAR(d_adj.close_time, 'HH24:MI')),
                        'branchsrc', (
                            SELECT
                                b.title 
                            FROM
                                r_adj
                            JOIN
                                branch AS b ON r_adj.transition_id = d_adj.transition_id AND b.id = r_adj.branch_id AND r_adj.rank = 1
                        ),
                        'stationsrc', (
                            SELECT
                                s.title 
                            FROM
                                r_adj
                            JOIN
                                station AS s ON r_adj.transition_id = d_adj.transition_id AND s.id = r_adj.station_id AND r_adj.rank = 1
                        ),
                        'branchdst', (
                            SELECT
                                b.title 
                            FROM
                                r_adj
                            JOIN
                                branch AS b ON r_adj.transition_id = d_adj.transition_id AND b.id = r_adj.branch_id AND r_adj.rank = 2
                        ),
\end{lstlisting}

\clearpage

\begin{lstlisting}[label=lst:get_chart_json_by_id:4,caption=Реализация функции получения схемы (часть 4),firstnumber=152]
                        'stationdst', (
                            SELECT
                                s.title 
                            FROM
                                r_adj
                            JOIN
                                station AS s ON r_adj.transition_id = d_adj.transition_id AND s.id = r_adj.station_id AND r_adj.rank = 2
                        )
                    )
                )
            FROM
                d_adj
            ), '[]'::JSONB)
        ) INTO result_json
    FROM
        chart AS c
    WHERE
        c.id = chart_id;

    RETURN result_json;
END;
$$ LANGUAGE plpgsql
\end{lstlisting}

\begin{lstlisting}[label=lst:get_route_json_by_id:1,caption=Реализация функции получения маршрута (часть 1),firstnumber=1]
CREATE OR REPLACE FUNCTION get_route_json_by_id(way_id INT)
RETURNS JSONB AS $$
DECLARE
    result_json JSONB;
BEGIN
    SELECT
        jsonb_build_object(
            'City', c.city,
            'ChartTitle', c.title,
            'Title', w.title,
            'Duration', w.duration,
            'RouteItems', COALESCE((
                WITH
                    wi AS (
                        SELECT
                            way_item.id,
                            way_item.way_id,
                            way_item.nexus,
                            way_item.step_nomer
                        FROM
                            way_item
                        WHERE
                            way_item.way_id = w.id
                    ),
                    ris AS (
                        SELECT
                            wis.way_item_id,
                            wis.station_id,
                            bs.branch_id,
                            s.title AS station_title,
                            b.title AS branch_title,
                            s.occupancy AS station_occupancy,
                            s.access AS station_access,
                            s.open_time AS station_open_time,
                            s.close_time AS station_close_time,
                            wi.step_nomer
\end{lstlisting}

\clearpage

\begin{lstlisting}[label=lst:get_route_json_by_id:2,caption=Реализация функции получения маршрута (часть 2),firstnumber=37]
                        FROM
                            wi
                        INNER JOIN
                            way_item_station AS wis ON wi.id = wis.way_item_id
                        INNER JOIN
                            station AS s ON s.id = wis.station_id
                        INNER JOIN
                            branch_station AS bs ON bs.station_id = s.id
                        INNER JOIN
                            branch AS b ON b.id = bs.branch_id
                    ),
                    rir AS (
                        SELECT
                            way_item_id,
                            branch_title,
                            from_station_title,
                            to_station_title,
                            railway_duration,
                            step_nomer
                        FROM (
                            SELECT
                                wir.way_item_id,
                                wi.nexus,
                                LAG(ris.branch_title) OVER (ORDER BY wi.step_nomer) AS branch_title,
                                LAG(ris.station_title) OVER (ORDER BY wi.step_nomer) AS from_station_title,
                                LEAD(ris.station_title) OVER (ORDER BY wi.step_nomer) AS to_station_title,
                                r.duration AS railway_duration,
                                wi.step_nomer
                            FROM
                                wi
                            LEFT JOIN
                                ris ON wi.id = ris.way_item_id
                            LEFT JOIN
                                way_item_railway AS wir ON wi.id = wir.way_item_id
                            LEFT JOIN
                                railway AS r ON wir.railway_id = r.id
                        )
                        WHERE
                            nexus = 'RAILWAY'::nexus_type
                    ),
                    rit AS (
                        SELECT
                            way_item_id,
                            from_station_title,
                            from_branch_title,
                            to_station_title,
                            to_branch_title,
                            transition_occupancy,
                            transition_access,
                            transition_duration,
                            transition_open_time,
                            transition_close_time,
                            step_nomer
\end{lstlisting}

\clearpage

\begin{lstlisting}[label=lst:get_route_json_by_id:3,caption=Реализация функции получения маршрута (часть 3),firstnumber=90]
                        FROM (
                            SELECT
                                wit.way_item_id,
                                wi.nexus,
                                LAG(ris.station_title) OVER (ORDER BY wi.step_nomer) AS from_station_title,
                                LAG(ris.branch_title) OVER (ORDER BY wi.step_nomer) AS from_branch_title,
                                LEAD(ris.station_title) OVER (ORDER BY wi.step_nomer) AS to_station_title,
                                LEAD(ris.branch_title) OVER (ORDER BY wi.step_nomer) AS to_branch_title,
                                t.occupancy AS transition_occupancy,
                                t.access AS transition_access,
                                t.duration AS transition_duration,
                                t.open_time AS transition_open_time,
                                t.close_time AS transition_close_time,
                                wi.step_nomer
                            FROM
                                wi
                            LEFT JOIN
                                ris ON wi.id = ris.way_item_id
                            LEFT JOIN
                                way_item_transition AS wit ON wi.id = wit.way_item_id
                            LEFT JOIN
                                transition AS t ON wit.transition_id = t.id
                            WHERE
                                wi.nexus != 'RAILWAY'::nexus_type
                        )
                        WHERE
                            nexus = 'TRANSITION'::nexus_type
                    ),
                    res AS (
                        SELECT
                            wi.step_nomer,
                            wi.nexus,
                            CASE wi.nexus
                                WHEN 'STATION'::nexus_type THEN ris.station_title
                                ELSE NULL
                            END AS station_title,
                            CASE wi.nexus
                                WHEN 'STATION'::nexus_type THEN ris.branch_title
                                WHEN 'RAILWAY'::nexus_type THEN rir.branch_title
                                ELSE NULL
                            END AS branch_title,
                            CASE wi.nexus
                                WHEN 'STATION'::nexus_type THEN ris.station_occupancy
                                ELSE NULL
                            END AS station_occupancy,
                            CASE wi.nexus
                                WHEN 'STATION'::nexus_type THEN ris.station_access
                                ELSE NULL
                            END AS station_access,
                            CASE wi.nexus
                                WHEN 'STATION'::nexus_type THEN ris.station_open_time
                                ELSE NULL
                            END AS station_open_time,
                            CASE wi.nexus
                                WHEN 'STATION'::nexus_type THEN ris.station_close_time
                                ELSE NULL
                            END AS station_close_time,
\end{lstlisting}

\clearpage

\begin{lstlisting}[label=lst:get_route_json_by_id:4,caption=Реализация функции получения маршрута (часть 4),firstnumber=147]
                            CASE wi.nexus
                                WHEN 'RAILWAY'::nexus_type THEN rir.from_station_title
                                WHEN 'TRANSITION'::nexus_type THEN rit.from_station_title
                                ELSE NULL
                            END AS from_station_title,
                            CASE wi.nexus
                                WHEN 'RAILWAY'::nexus_type THEN rir.to_station_title
                                WHEN 'TRANSITION'::nexus_type THEN rit.to_station_title
                                ELSE NULL
                            END AS to_station_title,
                            CASE wi.nexus
                                WHEN 'RAILWAY'::nexus_type THEN rir.railway_duration
                                ELSE NULL
                            END AS railway_duration,
                            CASE wi.nexus
                                WHEN 'TRANSITION'::nexus_type THEN rit.from_branch_title
                                ELSE NULL
                            END AS from_branch_title,
                            CASE wi.nexus
                                WHEN 'TRANSITION'::nexus_type THEN rit.to_branch_title
                                ELSE NULL
                            END AS to_branch_title,
                            CASE wi.nexus
                                WHEN 'TRANSITION'::nexus_type THEN rit.transition_occupancy
                                ELSE NULL
                            END AS transition_occupancy,
                            CASE wi.nexus
                                WHEN 'TRANSITION'::nexus_type THEN rit.transition_access
                                ELSE NULL
                            END AS transition_access,
                            CASE wi.nexus
                                WHEN 'TRANSITION'::nexus_type THEN rit.transition_duration
                                ELSE NULL
                            END AS transition_duration,
                            CASE wi.nexus
                                WHEN 'TRANSITION'::nexus_type THEN rit.transition_open_time
                                ELSE NULL
                            END AS transition_open_time,
                            CASE wi.nexus
                                WHEN 'TRANSITION'::nexus_type THEN rit.transition_close_time
                                ELSE NULL
                            END AS transition_close_time
                        FROM
                            wi
                        LEFT JOIN
                            ris ON wi.id = ris.way_item_id
                        LEFT JOIN
                            rir ON wi.id = rir.way_item_id
                        LEFT JOIN
                            rit ON wi.id = rit.way_item_id
                        ORDER BY
                            wi.step_nomer
                    )
                SELECT
                    jsonb_agg(
\end{lstlisting}

\clearpage

\begin{lstlisting}[label=lst:get_route_json_by_id:5,caption=Реализация функции получения маршрута (часть 5),firstnumber=202]
                        CASE res.nexus
                            WHEN 'STATION'::nexus_type THEN
                                jsonb_build_object(
                                    '$type', 'station',
                                    'Title', res.station_title,
                                    'BranchTitle', res.branch_title,
                                    'Occupancy', res.station_occupancy,
                                    'Access', res.station_access,
                                    'OpenTime', res.station_open_time,
                                    'CloseTime', res.station_close_time
                                )
                            WHEN 'RAILWAY'::nexus_type THEN
                                jsonb_build_object(
                                    '$type', 'railway',
                                    'BranchTitle', res.branch_title,
                                    'FromStationTitle', res.from_station_title,
                                    'ToStationTitle', res.to_station_title,
                                    'Duration', res.railway_duration
                                )
                            WHEN 'TRANSITION'::nexus_type THEN
                                jsonb_build_object(
                                    '$type', 'transition',
                                    'FromBranchTitle', res.from_branch_title,
                                    'FromStationTitle', res.from_station_title,
                                    'ToBranchTitle', res.to_branch_title,
                                    'ToStationTitle', res.to_station_title,
                                    'Occupancy', res.transition_occupancy,
                                    'Access', res.transition_access,
                                    'Duration', res.transition_duration,
                                    'OpenTime', res.transition_open_time,
                                    'CloseTime', res.transition_close_time
                                )
                            ELSE
                                jsonb_build_object('$type', 'unknown')
                        END
                    )
                FROM
                    res
            ), '[]'::JSONB)
        ) INTO result_json
    FROM
        way AS w
    INNER JOIN
        chart AS c ON w.chart_id = c.id
    WHERE
        w.id = way_id;


    RETURN result_json;
END;
$$ LANGUAGE plpgsql
\end{lstlisting}

\clearpage

\begin{lstlisting}[label=lst:color_convertors,caption=Реализация функций конвертации цвета,firstnumber=1]
CREATE OR REPLACE FUNCTION hex_color_to_int(hex_color TEXT)
RETURNS INT AS $$
DECLARE
    cleaned_hex TEXT;
    full_hex TEXT;
    color_int INT;
BEGIN
    cleaned_hex := UPPER(TRIM(BOTH '#' FROM hex_color));

    IF LENGTH(cleaned_hex) = 3 THEN
        full_hex := CONCAT(
            SUBSTR(cleaned_hex, 1, 1), SUBSTR(cleaned_hex, 1, 1),
            SUBSTR(cleaned_hex, 2, 1), SUBSTR(cleaned_hex, 2, 1),
            SUBSTR(cleaned_hex, 3, 1), SUBSTR(cleaned_hex, 3, 1)
        );
    ELSIF LENGTH(cleaned_hex) = 6 THEN
        full_hex := cleaned_hex;
    ELSE
        RAISE EXCEPTION 'Неверный формат hex-цвета: %', hex_color;
    END IF;

    EXECUTE 'SELECT ' || quote_literal('x' || full_hex) || '::bit(24)::int'
    INTO color_int;

    RETURN color_int;
END;
$$ LANGUAGE plpgsql IMMUTABLE;


CREATE OR REPLACE FUNCTION int_color_to_hex(int_color INT)
RETURNS TEXT AS $$
DECLARE
    hex_color TEXT;
BEGIN
    IF int_color IS NULL OR int_color < 0 OR int_color > 16777215 THEN
        RAISE EXCEPTION 'Цвет должен быть в диапазоне от 0 до 16777215, получено: %', int_color;
    END IF;

    hex_color := LPAD(TO_HEX(int_color), 6, '0');

    RETURN '#' || LOWER(hex_color);
END;
$$ LANGUAGE plpgsql IMMUTABLE;
\end{lstlisting}

\clearpage

\begin{lstlisting}[label=lst:nav_funcs,caption=Реализация навигационных функций,firstnumber=1]
CREATE OR REPLACE FUNCTION get_station_id_by_branch_id_station_title(
    p_branch_id INT,
    station_title TEXT
) RETURNS INT AS $$
DECLARE
    v_station_id INT;
BEGIN
    SELECT
        s.id INTO v_station_id
    FROM
        branch_station AS bs
    INNER JOIN
        station AS s ON bs.station_id = s.id
    WHERE
        bs.branch_id = p_branch_id AND
        s.title = station_title;
    
    RETURN v_station_id;
END;
$$ LANGUAGE plpgsql STABLE;


CREATE OR REPLACE FUNCTION get_station_id_by_branch_title_station_title(
    branch_title TEXT,
    station_title TEXT
) RETURNS INT AS $$
DECLARE
    v_station_id INT;
BEGIN
    SELECT
        s.id INTO v_station_id
    FROM
        branch_station AS bs
    INNER JOIN
        branch AS b ON bs.branch_id = b.id
    INNER JOIN
        station AS s ON bs.station_id = s.id
    WHERE
        b.title = branch_title AND
        s.title = station_title;
    
    RETURN v_station_id;
END;
$$ LANGUAGE plpgsql STABLE;
\end{lstlisting}

\clearpage

\begin{lstlisting}[label=lst:ck:tr:funcs:1,caption=Реализация проверяющий триггерных функций (часть 1),firstnumber=1]
-- Проверка корректного добавления связи Схема--Ветка
CREATE OR REPLACE FUNCTION ck_chart_branch_unique_ref() RETURNS TRIGGER
AS $$
BEGIN
    IF EXISTS (
        SELECT
            1
        FROM
            chart_branch
        WHERE
            chart_id <> NEW.chart_id AND branch_id = NEW.branch_id
    ) THEN
        RAISE EXCEPTION
            'Попытка связать ветку (%), уже связанную со схемой (%), со схемой (%).',
            branch_id, chart_id, NEW.chart_id;
    END IF;
    RETURN NEW;
END;
$$ LANGUAGE plpgsql;

-- Проверка корректного добавления связи Ветка--Станция
CREATE OR REPLACE FUNCTION ck_branch_station_unique_ref() RETURNS TRIGGER
AS $$
BEGIN
    IF EXISTS (
        SELECT
            1
        FROM
            branch_station
        WHERE
            branch_id <> NEW.branch_id AND station_id = NEW.station_id
    ) THEN
        RAISE EXCEPTION
            'Попытка связать станцию (%), уже связанную с веткой (%), с веткой (%).',
            station_id, branch_id, NEW.branch_id;
    END IF;
    RETURN NEW;
END;
$$ LANGUAGE plpgsql;

-- Проверка дублирования переездов между станциями
CREATE OR REPLACE FUNCTION ck_railway_unique_station_ref() RETURNS TRIGGER
AS $$
BEGIN
    IF EXISTS (
        SELECT
            1
        FROM
            railway
        WHERE
            from_id = NEW.to_id AND to_id = NEW.from_id
    ) THEN
        RAISE EXCEPTION
            'Дублирование переезда между станциями. Станции (%, %) уже соединены переездом: (%, %).',
            NEW.from_id, NEW.to_id, from_id, to_id;
    END IF;
    RETURN NEW;
END;
$$ LANGUAGE plpgsql;
\end{lstlisting}

\clearpage

\begin{lstlisting}[label=lst:ck:tr:funcs:2,caption=Реализация проверяющий триггерных функций (часть 2),firstnumber=61]
-- Проверка связывания станций переездом в рамках одной ветки
CREATE OR REPLACE FUNCTION ck_railway_same_branch_of_stations_ref() RETURNS TRIGGER
AS $$
DECLARE
    branch1_id INT;
    branch2_id INT;
BEGIN
    SELECT
        branch_id INTO branch1_id
    FROM
        branch_station AS bs
    WHERE
        NEW.from_id = bs.station_id
    LIMIT 1;

    SELECT
        branch_id INTO branch2_id
    FROM
        branch_station AS bs
    WHERE
        NEW.to_id = bs.station_id
    LIMIT 1;

    IF branch1_id IS NULL THEN
        RAISE EXCEPTION
            'Станция % не привязана ни к одной ветке',
            NEW.from_id;
    ELSIF branch2_id IS NULL THEN
        RAISE EXCEPTION
            'Станция % не привязана ни к одной ветке',
            NEW.to_id;
    ELSIF branch1_id <> branch2_id THEN
        RAISE EXCEPTION
            'Попытка соединения переездом станций (%, %), принадлежащие разным веткам: (%, %).',
            NEW.from_id, NEW.to_id, branch1_id, branch2_id;
    END IF;

    RETURN NEW;
END;
$$ LANGUAGE plpgsql;


-- Проверка связывания станций переходом в рамках разных веток одной и той же схемы
CREATE OR REPLACE FUNCTION ck_station_transition_different_branch_of_stations_ref() RETURNS TRIGGER
AS $$
DECLARE
    transition_links INT;
    station1_id INT;
    station2_id INT;
    branch1_id INT;
    branch2_id INT;
    chart1_id INT;
    chart2_id INT;
BEGIN
    -- Подсчет количества станций ведущих на один и тот же переход
    SELECT
        count(st.transition_id) INTO transition_links
    FROM
        station_transition AS st
    WHERE
        st.transition_id = NEW.transition_id;
\end{lstlisting}

\clearpage

\begin{lstlisting}[label=lst:ck:tr:funcs:3,caption=Реализация проверяющий триггерных функций (часть 3),firstnumber=123]
    -- Гарантия, что только две станции ведут на переход
    IF transition_links < 2 THEN
        RETURN NEW; -- переход еще не полностью установлен
    ELSIF transition_links > 2 THEN
        RAISE EXCEPTION
            'Попытка связать переход (%) с больше чем 2-мя станциями',
            NEW.transition_id;
    END IF;

    -- Первая станция уже имеется в новой записи NEW
    SELECT
        NEW.station_id INTO station1_id;

    -- Вторая станция избирается с учетом того, что всего станций 2 и проверкой
    -- на несовпадение с NEW.station_id
    SELECT
        st.station_id INTO station2_id
    FROM
        station_transition AS st
    WHERE
        st.transition_id = NEW.transition_id AND st.station_id <> NEW.station_id
    LIMIT 1;

    -- Вторая станция должна быть найдена
    IF station2_id IS NULL THEN
        RAISE EXCEPTION
            'Не удалось определить вторую станцию для перехода %',
            NEW.transition_id;
    END IF;

    -- Поиск ветки и схемы для первой станции
    SELECT
        bs.branch_id, cb.chart_id INTO branch1_id, chart1_id
    FROM
        branch_station bs
    LEFT JOIN
        chart_branch cb ON cb.branch_id = bs.branch_id
    WHERE
        bs.station_id = station1_id
    LIMIT 1;
    
    -- Поиск ветки и схемы для второй станции
    SELECT
        bs.branch_id, cb.chart_id INTO branch2_id, chart2_id
    FROM
        branch_station bs
    LEFT JOIN
        chart_branch cb ON cb.branch_id = bs.branch_id
    WHERE
        bs.station_id = station2_id
    LIMIT 1;

    IF branch1_id IS NULL THEN
        RAISE EXCEPTION
            'Станция (%) не привязана к ветке',
            station1_id;
    ELSIF branch2_id IS NULL THEN
        RAISE EXCEPTION
            'Станция (%) не привязана к ветке',
            station2_id;
    ELSIF chart1_id IS NULL THEN
        RAISE EXCEPTION
            'Ветка (%) не привязана к схеме',
            branch1_id;
\end{lstlisting}

\clearpage

\begin{lstlisting}[label=lst:ck:tr:funcs:4,caption=Реализация проверяющий триггерных функций (часть 4),firstnumber=187]
    ELSIF chart2_id IS NULL THEN
        RAISE EXCEPTION
            'Ветка (%) не привязана к схеме',
            branch2_id;
    ELSIF branch1_id = branch2_id THEN
        RAISE EXCEPTION
            'Попытка соединения переходом станций (%, %), принадлежащие одной и той же ветке: (%).',
            statoin1_id, station2_id, branch1_id;
    END IF;

    RETURN NEW;
END;
$$ LANGUAGE plpgsql;
\end{lstlisting}

\begin{lstlisting}[label=lst:del:tr:funcs,caption=Реализация удаляющих триггерных функций,firstnumber=1]
-- Удаление всех branch, которые были связаны с удаляемым chart
CREATE OR REPLACE FUNCTION delete_orphaned_branches_before_chart_delete()
RETURNS TRIGGER AS $$
BEGIN
    DELETE FROM branch
    WHERE id IN (
        SELECT cb.branch_id
        FROM chart_branch AS cb
        WHERE cb.chart_id = OLD.id
    );
    
    RETURN OLD;
END;
$$ LANGUAGE plpgsql;


-- Удаление всех Station, которые были связаны с удаляемым Branch
CREATE OR REPLACE FUNCTION delete_orphaned_stations_before_branch_delete()
RETURNS TRIGGER AS $$
BEGIN
    DELETE FROM station
    WHERE id IN (
        SELECT bs.station_id
        FROM branch_station AS bs
        WHERE bs.branch_id = OLD.id
    );
    
    RETURN OLD;
END;
$$ LANGUAGE plpgsql;


-- Удаление всех Transition, которые были связаны с удаляемым Station
CREATE OR REPLACE FUNCTION delete_orphaned_transitions_before_station_delete()
RETURNS TRIGGER AS $$
BEGIN
    DELETE FROM transition
    WHERE id IN (
        SELECT st.transition_id
        FROM station_transition AS st
        WHERE st.station_id = OLD.id
    );
    
    RETURN OLD;
END;
$$ LANGUAGE plpgsql;
\end{lstlisting}

\clearpage
